\begin{pastexam}[label={px:ncu-112-dsa-03}]{中央資工所 112 學年度 複選題第 3 題}
（複選；每題 5 分，答錯 1 題倒扣 1 分）

A leftist tree is a min tree satisfying that dist(RChild(i)) <= dist(LChild(i)), where dist(j)
denotes the number of edges on the shortest path from node j to a leaf node, RChild(i) and LChild(i)
denote the right child and left child of node i, respectively. Which of the following statements about
leftist trees are correct?
\begin{enumerate}[label=(\Alph*)]
\item The length of path to rightmost leaf is O(log n) for a leftist tree with n nodes.
\item Merging two leftist trees is to merge the right subtree of one tree with the other.
\item Delete min takes O(log n) time.
\item If the path to leftmost leaf has x nodes, then the leftist tree has at least $2^x-1$ nodes.
\end{enumerate}
\end{pastexam}

\begin{pastsoln}
答案是 (A)、(B)、(C)。

題目的 dist 數的是邊；資結書第~8.5~節的 $s(x)$ 數的是節點，空位 $s(\NIL) = 0$。 兩者差 $1$：
$\mathrm{dist}(x) = s(x) - 1$（空的子樹 dist 當成 $-1$），左傾的條件 $\mathrm{dist}(\text{右}) \le
\mathrm{dist}(\text{左})$ 和 $s(\text{右}) \le s(\text{左})$ 是同一件事。 下面用 $s$ 和資結書定理~8.2：
以 $x$ 為根的子樹至少有 $2^{s(x)} - 1$ 個節點，且從 $x$ 一路往右走到空位剛好經過 $s(x)$ 個節點。
\begin{itemize}
\item (A) 對。 由定理~8.2，$n \ge 2^{s(\text{根})} - 1$，最右邊的路徑有 $s(\text{根}) \le \lg(n + 1)$ 個
節點，長度 $O(\log n)$。
\item (B) 對。 合併兩棵最小左傾樹 $a$、$b$ 時，取根比較小的那棵（設為 $a$）當新的根，把 $a$ 的右子樹和
$b$ 遞迴合併，結果接回 $a$ 的右邊，必要時左右對調（資結書虛擬碼~8.9）。 敘述說的就是這個步驟，
「one tree」指根比較小的那棵。 遞迴只沿兩棵樹的右邊路徑走，由 (A) 是 $O(\log n)$。
\item (C) 對。 刪除最小值就是拿掉根，再把左、右兩棵子樹合併，一次合併 $O(\log n)$。
\item (D) 錯。 定理~8.2 講的是\textbf{最右邊}的路徑，最左邊的路徑可以很長。 反例：三個節點
$1, 2, 3$，$2$ 是 $1$ 的左子節點、$3$ 是 $2$ 的左子節點，都沒有右子節點。 每個節點的右邊是空的
（$s = 0$），左邊 $s \ge 0$，滿足左傾條件，也是最小樹。 最左邊的路徑有 $x = 3$ 個節點，整棵樹卻只有
$3 < 2^3 - 1 = 7$ 個節點。 一般地，一條往左的鏈有 $x$ 個節點、總數只有 $x$。
\end{itemize}
\end{pastsoln}
